% ===== UNIT 11: Stream 1: backbone architecture (3.2.1) =====
% Status: REWRITTEN  (change to REWRITTEN / REVIEWED when done; see CHECKLIST.md)
% Rules: style rewrite only. Keep every \cite, \ref, \label, number, equation and table value as is. Change the wording, not the content.
\subsubsection{Backbone architecture}
The backbone follows the TransReID architecture~\cite{he2021transreid}. It is a ViT-Base/16 with the SIE module and the JPM added. The SIE module encodes the 20 camera indices of VeRi-776 as learnable offsets, which are added to the patch token sequence. The JPM is a parallel branch. It randomly shuffles groups of patch tokens and uses a separate classification head. A BNNeck is applied on the 768-dimensional global CLS token, before the classification head~\cite{he2021transreid}.

The main difference from the original TransReID recipe is the backbone initialisation. We use \texttt{vit\_base\_\allowbreak patch16\_\allowbreak clip\_\allowbreak
224.openai} (OpenAI CLIP pretrained ViT-B/16 via
\texttt{timm}~\cite{radford2021clip}). It replaces the ImageNet-21k-only weights. CLIP pretraining gives richer visual and semantic features. They transfer better to fine-grained vehicle appearance.
